% ===== §10 The Design of Space and Perception =====
\section{The Design of Space and Perception}
\label{sec:space}

\noindent
Perception is formed, and it is formed somewhere. The convergence established that the perceptual frame is laid down by a history of engagement; but engagements have places, and the places are largely made. A person's everyday perception is trained, day after day, in built environments that were designed, and the design of those environments is therefore among the concrete means by which perception is cultivated or deadened. This section takes up the spatial dimension of aesthetic education, and it takes it up as a chapter in its own right because the theory of how space forms perception is deep and specific and cannot be compressed into a remark. The section sets out that theory in four registers, the phenomenological, the atmospheric, the design-theoretic, and the political-economic, each of which has studied how built space acts on perception; and it then states the paper's own contribution, that space trains or dulls the perception of those who dwell in it, that the home is the privileged relational field in which everyday perception is continually generated, and that the design of space carries a question of justice about whose perception is cared for and whose is starved.

\subsection{The phenomenology of built space}
\label{sub:space-phenomenology}

Architecture has its own phenomenology, and it reached, from the side of the made environment, the same embodiment the paper established from the side of the perceiving body. Pallasmaa's critique of the ocularcentrism of modern building is its sharpest statement. Architecture, he argued, has become an art of the eye alone, a retinal art of images fixed by the hurried glance of the camera, and in privileging vision it has starved the other senses through which a body actually inhabits a space, the touch of a handrail, the acoustic of a room, the temperature and smell and underfoot texture that make a place bodily habitable \citep{pallasmaa2005}. The eyes of the skin are the whole sensing body brought to bear on built space, and a building addressed only to the eye trains a diminished perception in those who live in it. Bachelard, earlier, had shown the inward face of the same relation: the spaces of a dwelling shape thought, memory, and reverie, so that the house is not a neutral container of a life but a formative structure of the perceiving and remembering that go on within it \citep{bachelard1994}. Norberg-Schulz gave the relation its name in the genius loci, the character or spirit of a place that a phenomenology of architecture must attend to, the specific way a place gathers and discloses a world to those who dwell there \citep{norbergschulz1980}. Across these the finding is one: built space is a formative condition of perception, addressing the sensing body, shaping the inward life, and disclosing a world with a specific character, so that to design a space is to design, in part, the perception of those who will inhabit it.

\subsection{Atmosphere and attunement}
\label{sub:space-atmosphere}

The aesthetics of atmosphere, introduced earlier as the aesthetic articulation of the relational locus of beauty (\xref{sub:gen-atmosphere}), returns here as the theory of what built space generates. What a designed space produces, on B\"{o}hme's account, is an atmosphere, a tuned space in which the qualities radiating from materials, proportions, and light meet the felt body of the one who enters, so that the space attunes its inhabitant to a mood before any object in it is regarded \citep{bohme2017}. The atmosphere is the space's primary aesthetic action on perception, the way it sets, beneath notice, the key in which everything within it will be perceived. Eastern spatial thought had refined a cognate insight into a discipline. The Japanese conception of \zh{間}, \emph{ma}, treats the interval, the pause, the charged emptiness between things, as itself a positive element of spatial and temporal composition, so that what is not built carries as much of the aesthetic work as what is; a space is attuned as much by its intervals and its held emptiness as by its filled and articulated parts. This will bear directly on one of the paper's models of practice, the leaving of emptiness (\xref{sec:typology}), for the designed interval is the built form of the field of contingency, a space made so that perception may generate within it rather than being directed to a prescribed object. Atmosphere and \emph{ma} together specify how space acts on perception below the level of the discrete object: it attunes, it sets a key, and it can do so as much by what it withholds as by what it presents.

\subsection{Affordance and the design of the everyday}
\label{sub:space-design}

Design theory proper supplies the most concrete register, and it does so through the concept the paper has already drawn on from ecological psychology. Gibson's affordances, the possibilities for action that a thing or place offers to a body whose capacities they answer to (\xref{sub:body-enactive}), became, in Norman's hands, the central concept of a theory of everyday design: a well-designed object communicates its affordances to the user, and where the affordance is not directly perceivable, a signifier is added to indicate it, so that design is the shaping of the perceived possibilities of action that a made thing offers \citep{gibson1979,norman2013}. The bearing on perception is exact and everyday. The objects and spaces of ordinary life, in offering or failing to offer perceivable affordances, train the perception and action of those who use them; a well-made handle teaches the hand, a coherent room teaches the movement through it, and the daily traffic with designed things is a continual, unnoticed schooling of perception in what the made world affords. Norman's title names the paper's own domain, the design of everyday things, and his analysis supplies the fine grain of how designed everyday objects act on the perceiving body that uses them. What his user-centred design largely set aside, the aesthetic and not merely the functional dimension of this shaping, is what the present chapter restores to it: the affordances a space offers are not only for use but for perception, and a space may afford, or fail to afford, the finely generated aesthetic perception of an ordinary life.

\subsection{The production of perceptual space}
\label{sub:space-politicaleconomy}

The formation of perception by space is not innocent of power, and the political economy of the senses established earlier has a spatial extension that this register supplies. Space, Lefebvre argued, is produced: it is not a neutral container in which social life happens but a social product, produced by and reproducing the relations of a mode of production, so that the organisation of space is among the means by which a social order maintains itself \citep{lefebvre1991}. Joined to the claim that the senses are historically produced (\xref{sub:val-politicaleconomy}), this yields the proposition the paper needs here: the spaces in which everyday perception is formed are themselves produced under a mode of production, and they form perception in its interest. The retail environment engineered to capture attention and direct it to purchase, the domestic space arranged around the screens that resell the attention they absorb, the built environments that afford consumption richly and generative perception hardly at all, are the spatial production of a perception fitted to an economy. The habituation and numbing diagnosed at the paper's outset are trained, in part, by spaces built to train them. This is the spatial face of the commodified training of perception marked as the second spine's corruption (\xref{sub:plas-redline}): perception is reorganised toward consumption not only by what it is shown but by the very spaces it is shown it in, and the design of space is therefore one of the principal sites at which the cultivation and the corruption of everyday perception are decided.

\subsection{The home as a field of perceptual generation}
\label{sub:space-home}

The paper's own contribution gathers these registers to a claim about the home, which is, for the everyday perception this paper cares about most, the decisive designed space. The home is where the greatest part of everyday aesthetic life is generated, the space whose light and sound and texture and arrangement make up the sensory substance of ordinary days, and it is, unlike the retail and the public spaces produced by a distant power, a space its inhabitants themselves in large part make and arrange. This gives the home a double character in the paper's argument. It is, first, the built condition under which a shared everyday perception is trained: the arrangement of a home, its affordances and atmospheres and intervals, forms the perception of those who dwell in it, and a home may be arranged to afford the finely generated perception of ordinary life or to starve it, to attune its inhabitants to their shared days or to distract them from these into the resold attention of the screen. It is, second and distinctively, a space of relational making, since a shared home is arranged by those who share it, and its arrangement is itself one of the practices by which they cultivate between them the perception of their common life. The home is thus not only the space in which relational aesthetic education occurs but one of its media, a designed field that its inhabitants continually remake and that continually reforms their shared perception in turn, in the spiral by which the individuals and the relation form each other. This connects the chapter forward to the models of practice, where the arrangement of a shared space by benevolent signs becomes one of the concrete forms of everyday aesthetic generation (\xref{sec:typology}), and it carries the chapter's justice to the threshold of the paper's account of justice proper. For if perception is formed by designed space, and if some are given spaces that afford the cultivation of perception while others are given spaces built to starve it and resell their attention, then the design of space distributes perceptual capacity unequally, and does so along the lines of who can command the making of their own space. Whose home affords the generation of a shared aesthetic life, and whose does not, is a question of justice, and it is to justice that the paper now turns (\xref{sec:distribution}).
